% ============================================================
%  Impacto da IA em Repositórios GitHub — Lab04
%  Compilar no Overleaf com pdfLaTeX
% ============================================================
\documentclass[12pt,a4paper,twocolumn]{article}

\usepackage[utf8]{inputenc}
\usepackage[T1]{fontenc}
\usepackage[brazil]{babel}
\usepackage{geometry}
\usepackage{graphicx}
\usepackage{booktabs}
\usepackage{tabularx}
\usepackage{pgfplots}
\usepackage{tikz}
\usepackage{xcolor}
\usepackage{hyperref}
\usepackage{caption}
\usepackage{amsmath}
\usepackage{microtype}
\usepackage{enumitem}
\usepackage{float}
\usepackage{parskip}
\usepackage{titlesec}
\usepackage{fancyhdr}
\usepackage{authblk}
\usepackage{mdframed}

\raggedbottom

\pgfplotsset{compat=1.18}

\geometry{a4paper, top=2.4cm, bottom=2.4cm, left=1.8cm, right=1.8cm, columnsep=0.7cm}

% ── Cores ──────────────────────────────────────────────────
\definecolor{primblue}{RGB}{30,100,200}
\definecolor{acyan}{RGB}{0,180,210}
\definecolor{agreen}{RGB}{0,170,120}
\definecolor{ared}{RGB}{210,55,70}
\definecolor{aamber}{RGB}{200,140,10}
\definecolor{darkgray}{RGB}{80,100,120}
\definecolor{lightpanel}{RGB}{248,250,253}
\definecolor{borderpanel}{RGB}{210,220,235}

\hypersetup{colorlinks=true,linkcolor=primblue,citecolor=primblue,urlcolor=acyan}

% ── Seções ─────────────────────────────────────────────────
\titleformat{\section}{\large\bfseries\color{primblue}}{\thesection.}{0.5em}{}[\vspace{-4pt}\rule{\columnwidth}{0.4pt}]
\titleformat{\subsection}{\normalsize\bfseries}{\thesubsection.}{0.5em}{}
\titlespacing{\section}{0pt}{10pt}{4pt}
\titlespacing{\subsection}{0pt}{6pt}{2pt}

% ── Cabeçalho ──────────────────────────────────────────────
\pagestyle{fancy}\fancyhf{}
\fancyhead[L]{\small\color{darkgray}\textit{Lab04 — Impacto da IA em Repositórios GitHub}}
\fancyhead[R]{\small\color{darkgray}2025}
\fancyfoot[C]{\small\thepage}
\renewcommand{\headrulewidth}{0.4pt}
\renewcommand{\headrule}{\hbox to\headwidth{\color{acyan}\leaders\hrule height\headrulewidth\hfill}}

% ── Box de hipóteses ───────────────────────────────────────
\newmdenv[backgroundcolor=lightpanel,linecolor=borderpanel,linewidth=0.5pt,
          innerleftmargin=6pt,innerrightmargin=6pt,
          innertopmargin=5pt,innerbottommargin=5pt,
          skipabove=4pt,skipbelow=4pt, nobreak=true]{hypbox}

% ── Pgfplots ───────────────────────────────────────────────
\pgfplotsset{
  pairbar/.style={
    width=0.95\columnwidth, height=5.2cm,
    ybar, bar width=14pt, enlarge x limits=0.35,
    symbolic x coords={Mediana, Media}, xtick=data,
    ymin=0, grid=major, grid style={dotted,gray!40},
    legend style={font=\tiny,at={(0.5,1.02)},anchor=south,legend columns=-1,draw=none,fill=none},
    tick label style={font=\tiny}, label style={font=\tiny},
    title style={font=\small\bfseries,at={(0.5,1.15)},anchor=south},
    cycle list={{fill=acyan!70,draw=none},{fill=agreen!70,draw=none}},
  }
}

\begin{document}

% ── Título (onecolumn) ─────────────────────────────────────
\twocolumn[{%
\begin{@twocolumnfalse}
\vspace*{0.2cm}
{\centering
  {\LARGE\bfseries Impacto de Ferramentas de Inteligência Artificial\\[4pt]
   em Repositórios GitHub: Uma Análise Empírica\par}
  \vspace{0.5em}
  {\large Andre Lazarini e João Pedro Guimarães\par}
  \vspace{0.15em}
  {\small Laboratório de Experimentação em Engenharia de Software\par}
\par}

\vspace{0.5em}\noindent\rule{\linewidth}{0.5pt}

\begin{center}\textbf{Resumo}\end{center}
\small
Este trabalho apresenta uma análise empírica de 105 repositórios Python públicos do
GitHub classificados por nível de evidência de uso de IA. Para cada repositório,
comparamos métricas de \textit{issues}, retrabalho corretivo e verbosidade de
\textit{commits} nos períodos anterior e posterior ao primeiro sinal de adoção de IA.
Os resultados indicam: \textbf{(RQ1)} redução na taxa de \textit{issues} e no tempo
de resolução — contrário à hipótese alternativa; \textbf{(RQ2)} aumento modesto em
\textit{fixes} e \textit{reverts} — em linha com a hipótese alternativa;
\textbf{(RQ3)} \textit{commits} maiores pós-IA, mas sem melhora no Maintainability
Index — suporte parcial.

\noindent\rule{\linewidth}{0.5pt}\vspace{0.5em}
\end{@twocolumnfalse}
}]

% ═══════════════════════════════════════════════════════════
\section{Introdução}
% ═══════════════════════════════════════════════════════════

Ferramentas de IA generativa — GitHub Copilot, Claude, ChatGPT — integraram-se
ao fluxo de desenvolvimento de inúmeros projetos de código aberto. Estudos prévios
investigaram a qualidade do código gerado por IA~\cite{dakhel2023copilot} e a
produtividade de desenvolvedores~\cite{ziegler2022copilot}, mas há lacuna em
análises empíricas de larga escala que comparem métricas \textit{antes} e
\textit{depois} da introdução de sinais de IA nos repositórios.

Este trabalho aborda três questões de pesquisa:

\begin{itemize}[leftmargin=1.2em,itemsep=2pt,topsep=2pt]
  \item \textbf{RQ1:} O uso de IA impacta a taxa de abertura de \textit{issues}
    e o tempo médio de resolução em projetos de software?
  \item \textbf{RQ2:} \textit{Commits} assistidos por IA exigem mais retrabalho
    corretivo?
  \item \textbf{RQ3:} \textit{Commits} assistidos por IA são mais verbosos
    (maiores)?
\end{itemize}

% ═══════════════════════════════════════════════════════════
\section{Trabalhos Relacionados}
% ═══════════════════════════════════════════════════════════

Dakhel et al.~\cite{dakhel2023copilot} identificaram limitações do Copilot em
lógica complexa. Ziegler et al.~\cite{ziegler2022copilot} mediram a aceitação
de sugestões por profissionais. Nguyen e Nadi~\cite{nguyen2022copilot} avaliaram
corretude funcional de código gerado por IA.

Diferentemente desses trabalhos, nossa abordagem é \textit{observacional} e baseada
em mineração de repositórios reais, comparando métricas temporais antes e depois da
introdução de sinais de IA, sem experimentos controlados.

% ═══════════════════════════════════════════════════════════
\section{Metodologia}
% ═══════════════════════════════════════════════════════════

\subsection{Coleta e Seleção}

Foram selecionados 105 repositórios Python públicos do GitHub com pelo menos 1.000
estrelas e ao menos um sinal detectado de uso de IA. A coleta foi feita via API do
GitHub e análise estática de código para o Maintainability Index (MI).

\subsection{Sinais de IA e Agrupamento}

Quatro tipos de evidência foram investigados: (1)~co-autoria em \textit{commits};
(2)~menção ao Claude; (3)~menção ao Copilot; (4)~arquivo de configuração de IA.
O \textit{AI Score} = total de evidências encontradas. Grupos: \textbf{Baixo}
(2--3), \textbf{Médio} (4--6), \textbf{Alto} ($\geq 7$).

\subsection{Métricas por Questão de Pesquisa}

\begin{itemize}[leftmargin=1.2em, itemsep=2pt, parsep=0pt, partopsep=0pt, topsep=2pt]
  \item \textbf{RQ1 — Issues:} M1 = taxa de \textit{issues}/mês; M2 = tempo
    médio de vida das \textit{issues} (horas).
  \item \textbf{RQ2 — Retrabalho:} M1 = proporção de \textit{commits}
    corretivos (\textit{fixes}); M2 = taxa de \textit{commits} de revert.
  \item \textbf{RQ3 — Verbosidade:} M1 = linhas alteradas por \textit{commit};
    M2 = Maintainability Index (MI)~\cite{oman1992maintainability}.
\end{itemize}

\subsection{Hipóteses}

Para cada RQ definimos hipótese nula (H0) e alternativa (H1):

\begin{hypbox}
\footnotesize
\textbf{RQ1 — H0:} sem diferença em taxa de \textit{issues} nem em tempo de resolução
pré/pós IA.\\
\textbf{RQ1 — H1:} IA → maior taxa de \textit{issues} e maior tempo de resolução.
\end{hypbox}

\begin{hypbox}
\footnotesize
\textbf{RQ2 — H0:} proporção de \textit{fixes} e reverts iguais com e sem IA.\\
\textbf{RQ2 — H1:} IA → maior proporção de \textit{commits} corretivos e mais reverts.
\end{hypbox}

\begin{hypbox}
\footnotesize
\textbf{RQ3 — H0:} sem diferença no tamanho de \textit{commits} nem no MI.\\
\textbf{RQ3 — H1:} IA → \textit{commits} maiores (mais linhas) e maior complexidade.
\end{hypbox}

\subsection{Análise}

Para cada repositório, identificamos o primeiro sinal de IA (marco temporal) e
calculamos as métricas nos períodos anterior e posterior. Adotamos a \textbf{mediana}
como medida principal para reduzir o efeito de \textit{outliers}.

% ═══════════════════════════════════════════════════════════
\section{Caracterização do Dataset}
% ═══════════════════════════════════════════════════════════

A Tabela~\ref{tab:dataset} resume a amostra e a Tabela~\ref{tab:grupos} detalha os
grupos de adoção.

\begin{table}[t]
\centering\caption{Caracterização geral.}\label{tab:dataset}
\small
\begin{tabular}{lc}
\toprule
\textbf{Atributo} & \textbf{Valor} \\
\midrule
Repositórios       & 105 \\
Mediana de estrelas & 17.169 \\
Mediana de idade   & 9,3 anos \\
Mediana AI Score   & 3 \\
\bottomrule
\end{tabular}
\end{table}

\begin{table}[t]
\centering\caption{Grupos de adoção de IA.}\label{tab:grupos}
\small\setlength{\tabcolsep}{4pt}
\begin{tabular}{lrrrr}
\toprule
\textbf{Grupo} & \textbf{n} & \textbf{Med.~$\bigstar$} & \textbf{Med. Idade} & \textbf{Med. Score} \\
\midrule
Baixo (2--3) & 55 & 14.662 & 9,50 a & 3,0 \\
Médio (4--6) & 32 & 18.905 & 9,00 a & 6,0 \\
Alto ($\geq 7$) & 18 & 30.102 & 9,05 a & 8,5 \\
\bottomrule
\end{tabular}
\end{table}

\begin{figure}[t]
\centering
\begin{tikzpicture}
\begin{axis}[
  width=0.95\columnwidth, height=4cm,
  ybar, bar width=20pt,
  symbolic x coords={Baixo (2-3),Medio (4-6),Alto (7+)},
  xtick=data,
  xticklabel style={font=\footnotesize},
  ymin=0, ymax=65, nodes near coords, nodes near coords style={font=\footnotesize},
  grid=major, grid style={dotted,gray!40},
  title={\footnotesize Distribuição por grupo de adoção de IA},
  title style={font=\small\bfseries,at={(0.5,1.06)},anchor=south},
  enlarge x limits=0.2,
  ylabel={\footnotesize Repositórios}, ylabel style={font=\footnotesize},
]
\addplot[fill=aamber!80,draw=none] coordinates {
  (Baixo (2-3),55) (Medio (4-6),32) (Alto (7+),18)
};
\end{axis}
\end{tikzpicture}
\caption{Repositórios por grupo.}
\label{fig:grupos}
\end{figure}

% ═══════════════════════════════════════════════════════════
\section{Resultados}
% ═══════════════════════════════════════════════════════════

\subsection{RQ1 — Taxa de Issues e Tempo de Resolução}

\begin{table}[t]
\centering\caption{RQ1 — medianas antes e depois.}\label{tab:rq1}
\small\setlength{\tabcolsep}{4pt}
\begin{tabular}{lrrr}
\toprule
\textbf{Métrica} & \textbf{Antes} & \textbf{Depois} & \textbf{$\Delta$} \\
\midrule
Taxa issues/mês & 66,50 & 18,83 & \textcolor{ared}{$-$71,7\%} \\
Tempo resolução (h) & 9.898,01 & 5.645,14 & \textcolor{ared}{$-$42,9\%} \\
\bottomrule
\end{tabular}
\end{table}

\begin{figure}[t]
\centering
\begin{tikzpicture}
\begin{axis}[
  pairbar,
  ylabel={\footnotesize Taxa issues/mês},
  title={\footnotesize RQ1 — M1: Taxa de issues/mês},
  ymax=280, legend entries={Antes,Depois},
]
\addplot coordinates { (Mediana,66.5) (Media,234.31) };
\addplot coordinates { (Mediana,18.83) (Media,75.39) };
\end{axis}
\end{tikzpicture}
\caption{RQ1 — M1: taxa de issues/mês antes e depois.}
\label{fig:rq1m1}
\end{figure}

\begin{figure}[t]
\centering
\begin{tikzpicture}
\begin{axis}[
  pairbar,
  ylabel={\footnotesize Tempo de resolução (h)},
  title={\footnotesize RQ1 — M2: Tempo médio de resolução},
  ymax=14000, legend entries={Antes,Depois},
]
\addplot coordinates { (Mediana,9898) (Media,12338) };
\addplot coordinates { (Mediana,5645) (Media,9290) };
\end{axis}
\end{tikzpicture}
\caption{RQ1 — M2: tempo médio de resolução (horas) antes e depois.}
\label{fig:rq1m2}
\end{figure}

\textbf{Veredicto — RQ1:} H1 \textbf{rejeitada}. Ambas as métricas diminuem
pós-IA (taxa $-71{,}7\%$; tempo $-42{,}9\%$), contrariando a hipótese de que
a IA aumentaria o fluxo de \textit{issues}.

\subsection{RQ2 — Retrabalho Corretivo}

\begin{table}[t]
\centering\caption{RQ2 — medianas antes e depois.}\label{tab:rq2}
\small\setlength{\tabcolsep}{4pt}
\begin{tabular}{lrrr}
\toprule
\textbf{Métrica} & \textbf{Antes} & \textbf{Depois} & \textbf{$\Delta$} \\
\midrule
\% Fixes  & 32,12 & 36,07 & \textcolor{agreen}{$+$12,3\%} \\
\% Reverts & 1,10 &  1,24 & \textcolor{agreen}{$+$12,7\%} \\
\bottomrule
\end{tabular}
\end{table}

\begin{figure}[t]
\centering
\begin{tikzpicture}
\begin{axis}[
  pairbar,
  ylabel={\footnotesize \% Fixes (mediana)},
  title={\footnotesize RQ2 — M1: Proporção de commits corretivos},
  ymax=44, legend entries={Antes,Depois},
]
\addplot coordinates { (Mediana,32.12) (Media,33.45) };
\addplot coordinates { (Mediana,36.07) (Media,38.76) };
\end{axis}
\end{tikzpicture}
\caption{RQ2 — M1: percentual de \textit{commits} de correção antes e depois.}
\label{fig:rq2m1}
\end{figure}

\textbf{Veredicto — RQ2:} H1 \textbf{parcialmente suportada}. Tanto
\textit{fixes} ($+12{,}3\%$) quanto \textit{reverts} ($+12{,}7\%$) aumentam
após a adoção de IA, em linha com a hipótese alternativa.

\subsection{RQ3 — Verbosidade dos Commits}

\begin{table}[t]
\centering\caption{RQ3 — medianas antes e depois.}\label{tab:rq3}
\small\setlength{\tabcolsep}{4pt}
\begin{tabular}{lrrr}
\toprule
\textbf{Métrica} & \textbf{Antes} & \textbf{Depois} & \textbf{$\Delta$} \\
\midrule
Linhas/commit    & 193,79 & 217,49 & \textcolor{agreen}{$+$12,2\%} \\
MI médio (méd.) & 63,43  & 60,90  & \textcolor{ared}{$-$2,53} \\
MI (mediana)    & 60,29  & 58,50  & \textcolor{ared}{$-$1,80} \\
\bottomrule
\end{tabular}
\end{table}

\begin{figure}[t]
\centering
\begin{tikzpicture}
\begin{axis}[
  pairbar,
  ylabel={\footnotesize Linhas alteradas/commit},
  title={\footnotesize RQ3 — M1: Tamanho médio do commit},
  ymax=700, legend entries={Antes,Depois},
]
\addplot coordinates { (Mediana,193.79) (Media,573.01) };
\addplot coordinates { (Mediana,217.49) (Media,591.69) };
\end{axis}
\end{tikzpicture}
\caption{RQ3 — M1: linhas alteradas por \textit{commit} antes e depois.}
\label{fig:rq3m1}
\end{figure}

\begin{figure}[t]
\centering
\begin{tikzpicture}
\begin{axis}[
  pairbar,
  ylabel={\footnotesize Maintainability Index},
  title={\footnotesize RQ3 — M2: Maintainability Index},
  ymin=54, ymax=68, legend entries={Antes,Depois},
]
\addplot coordinates { (Mediana,60.29) (Media,63.43) };
\addplot coordinates { (Mediana,58.50) (Media,60.90) };
\end{axis}
\end{tikzpicture}
\caption{RQ3 — M2: Maintainability Index antes e depois.}
\label{fig:rq3m2}
\end{figure}

\textbf{Veredicto — RQ3:} H1 \textbf{parcialmente suportada}. Linhas por
\textit{commit} aumentam $+12{,}2\%$ (suporta H1 para tamanho), mas o MI cai
ligeiramente (não suporta H1 para complexidade — código não fica mais simples).

% ═══════════════════════════════════════════════════════════
\section{Discussão}
% ═══════════════════════════════════════════════════════════

\textbf{RQ1 — achado contraintuitivo.} A hipótese alternativa previa mais
\textit{issues} e mais tempo de resolução após IA. Os dados mostram o oposto.
Possíveis explicações: (i)~a IA aumenta eficiência de triagem e resolução;
(ii)~o código gerado por IA apresenta menos defeitos observáveis;
(iii)~o período pós-sinal é naturalmente mais curto, acumulando menos eventos;
(iv)~projetos mais maduros tendem a ter atividade decrescente
independentemente da IA.

\textbf{RQ2 — retrabalho leve, mas consistente.} Ambas as métricas sobem
de forma modesta (~12\%). O aumento de \textit{fixes} e \textit{reverts}
sugere que código assistido por IA pode exigir ajustes posteriores, mas
o efeito não é dramático.

\textbf{RQ3 — commits maiores, sem ganho de qualidade estrutural.} O
aumento de linhas por \textit{commit} pode refletir que a IA gera blocos
de código maiores por interação. A queda no MI, embora pequena (~3\%),
indica que a qualidade estrutural não melhora com a IA — o código gerado
pode ser ligeiramente mais difícil de manter.

% ═══════════════════════════════════════════════════════════
\section{Ameaças à Validade}
% ═══════════════════════════════════════════════════════════

\textbf{Interna:} estudo observacional sem grupo de controle; causalidade não pode
ser inferida. Variáveis de confusão (maturidade, mudança de equipe) não controladas.

\textbf{Construto:} AI Score mede evidência, não intensidade de uso. MI é sensível
à linguagem e às ferramentas de análise.

\textbf{Externa:} amostra restrita a Python com $\geq 1{.}000$ estrelas; pode não
generalizar para projetos menores ou outras linguagens.

\textbf{Conclusão:} mediana mitiga \textit{outliers}, mas repositórios com
comportamento atípico ainda podem influenciar os resultados.

% ═══════════════════════════════════════════════════════════
\section{Conclusão}
% ═══════════════════════════════════════════════════════════

Este trabalho analisou empiricamente 105 repositórios Python quanto ao impacto
de ferramentas de IA. Os resultados são:

\begin{itemize}[leftmargin=1.2em,itemsep=2pt,topsep=2pt]
  \item \textbf{RQ1:} taxa de \textit{issues} e tempo de resolução
    \textbf{diminuem} pós-IA — H1 rejeitada.
  \item \textbf{RQ2:} \textit{fixes} e \textit{reverts} \textbf{aumentam
    moderadamente} — H1 parcialmente suportada.
  \item \textbf{RQ3:} \textit{commits} \textbf{ficam maiores}, mas MI
    \textbf{cai levemente} — H1 parcialmente suportada.
\end{itemize}

Trabalhos futuros devem incluir grupos de controle, análise qualitativa de
\textit{commits} assistidos por IA e replicação em outras linguagens.

% ── Referências ───────────────────────────────────────────
\bibliographystyle{plain}
\begin{thebibliography}{9}

\bibitem{dakhel2023copilot}
  A.~M. Dakhel et al.,
  ``GitHub Copilot AI pair programmer: Asset or Liability?''
  \textit{J.\ Syst.\ Softw.}, 203:111806, 2023.

\bibitem{ziegler2022copilot}
  A.~Ziegler et al.,
  ``Productivity Assessment of Neural Code Completion.''
  In: \textit{Proc.\ MAPS}, 2022.

\bibitem{nguyen2022copilot}
  N.~Nguyen e S.~Nadi,
  ``An empirical evaluation of GitHub Copilot's code suggestions.''
  In: \textit{Proc.\ MSR}, pp.\ 1--5, 2022.

\bibitem{oman1992maintainability}
  P.~Oman e J.~Hagemeister,
  ``Metrics for assessing a software system's maintainability.''
  In: \textit{Proc.\ ICSM}, pp.\ 337--344, 1992.

\bibitem{kalliamvakou2014github}
  E.~Kalliamvakou et al.,
  ``The promises and perils of mining GitHub.''
  In: \textit{Proc.\ MSR}, pp.\ 92--101, 2014.

\end{thebibliography}

\end{document}